% Einheit 5 von 5 – Ortskurve und Kurvenvergleich (bank/funktionsscharen-und-ortskurven/e5.jsonl, zusammenbau v0.1)
%% TODO Zeitmarke Einheit 5: Marken-Zeile nicht lesbar
%% TODO Zweigzeile Teil 1: was hier gelernt wird (Satz in Schülersprache aus dem Einheitstitel) – Einheit 5
\einheitenkopf[e5]{Einheit 5 von 5 · Ortskurve und Kurvenvergleich}
\zweigzeile{Abitur LK}
\setcounter{aufgabe}{17}

%% TODO Titel: Ich-kann-Satz fehlt in der Bank (Platzhalter: Kettenname) – Nr. 18
%% TODO Anweisung: ein Satz über der ersten Teilaufgabe fehlt in der Bank – Nr. 18
\begin{aufgabe}{Ortskurve und Kurvenvergleich}
%% TODO \rechenplatz{n}: Schrittzahl des Grundfalls fehlt in der Bank (nur bei fester Schreibform, 2.3 b)
\begin{teile}
\teil Schar $f_a(x) = x^2 - 2a x$, Scheitel $S(a \,|\, -a^2)$. Trägt der Punkt den Parameter oder ist er parameterfrei? \\ \kreuz{trägt den Parameter} \kreuz{parameterfrei}
\teil Bestimme die Ortskurve der Scheitelpunkte der Schar $f_a(x) = x^2 - 2a x$. $y =$ \_\_
\teil Die Tiefpunkte der Schar $f_a(x) = x^2 - 2a x + 3$ liegen auf der Kurve $y = 3 - x^2$. Prüfe das für $a = 2$.
\teil Die Hochpunkte einer Schar liegen auf einer Geraden; zwei davon sind $H_1(2 | 5)$ und $H_2(4 | 9)$. Bestimme die Steigung und die Gleichung dieser Ortsgeraden. $m =$ \_\_, $y =$ \_\_
\teil Gegeben ist $f_a(x) = x^3 \cdot e^{-a x}$ mit $a > 0$; der Hochpunkt liegt bei $x = \frac{3}{a}$. Weise nach, dass er für jedes $a$ auf dem Graphen von $y = \frac{x^3}{e^3}$ liegt. \hfill (Abitur 2017 LK)
\teil Gegeben ist $p_k(x) = -(x - k)^2 + 2k$. Berechne den Abstand der Hochpunkte zweier benachbarter Parabeln $p_k$ und $p_{k+1}$. \leerfeld
\teil Gegeben sind $g_b(x) = \frac{1}{4}x \cdot (x - b) \cdot e^{x}$ mit $b \ge 0$ und $f(x) = \frac{1}{4}x \cdot (5 - x) \cdot e^{x}$. Für welches $b$ ist der Graph von $g_b$ das Spiegelbild des Graphen von $f$ an der x-Achse? \hfill (Abitur 2020 GK) \\ $b =$ \_\_
\teil Bestimme die gemeinsamen Punkte der Graphen von $f_k(x) = k x^2$ mit $k \neq 0$ und ihrer Ableitung $f_k'$.
\teil Gegeben ist $f_a(x) = \frac{3}{2}x - \frac{a^2}{2}x^3$ mit $a > 0$. Der Graph von $f_a$ entsteht aus dem Graphen von $f_1$ durch Streckung vom Ursprung aus mit dem Faktor $\frac{1}{a}$ in x- und in y-Richtung. Begründe ohne Berechnung der allgemeinen Extrempunkte, dass alle Extrempunkte der Schar auf der Geraden $y = x$ liegen. \hfill (Abitur 2022 LK)
\end{teile}
\end{aufgabe}

%% TODO Titel: Ich-kann-Satz fehlt in der Bank (Platzhalter: Kettenname) – Nr. 19
%% TODO Anweisung: ein Satz über der ersten Teilaufgabe fehlt in der Bank – Nr. 19
\begin{aufgabe}{Ortskurve und Kurvenvergleich – Fehler finden, Begründen}
\begin{teile}
\teil Die Tiefpunkte von $f_a(x) = x^2 - 4a x$ sind $T(2a \,|\, -4a^2)$. Paul schreibt: „Mit $x = 2a$ ersetze ich $a$ durch $x$: $y = -4x^2$.“ Finde den Fehler und bestimme die Ortskurve richtig.
\teil Die Extrempunkte einer Schar sind $E(a \,|\, 2a^2)$. Begründe, warum die Gleichung $y = 2x^2$, die man beim Eliminieren erhält, alle Extrempunkte zugleich beschreibt.
\end{teile}
\end{aufgabe}
